% Options for packages loaded elsewhere
\PassOptionsToPackage{unicode}{hyperref}
\PassOptionsToPackage{hyphens}{url}
\documentclass[
]{article}
\usepackage{xcolor}
\usepackage[margin=1in]{geometry}
\usepackage{amsmath,amssymb}
\setcounter{secnumdepth}{-\maxdimen} % remove section numbering
\usepackage{iftex}
\ifPDFTeX
  \usepackage[T1]{fontenc}
  \usepackage[utf8]{inputenc}
  \usepackage{textcomp} % provide euro and other symbols
\else % if luatex or xetex
  \usepackage{unicode-math} % this also loads fontspec
  \defaultfontfeatures{Scale=MatchLowercase}
  \defaultfontfeatures[\rmfamily]{Ligatures=TeX,Scale=1}
\fi
\usepackage{lmodern}
\ifPDFTeX\else
  % xetex/luatex font selection
\fi
% Use upquote if available, for straight quotes in verbatim environments
\IfFileExists{upquote.sty}{\usepackage{upquote}}{}
\IfFileExists{microtype.sty}{% use microtype if available
  \usepackage[]{microtype}
  \UseMicrotypeSet[protrusion]{basicmath} % disable protrusion for tt fonts
}{}
\makeatletter
\@ifundefined{KOMAClassName}{% if non-KOMA class
  \IfFileExists{parskip.sty}{%
    \usepackage{parskip}
  }{% else
    \setlength{\parindent}{0pt}
    \setlength{\parskip}{6pt plus 2pt minus 1pt}}
}{% if KOMA class
  \KOMAoptions{parskip=half}}
\makeatother
\usepackage{color}
\usepackage{fancyvrb}
\newcommand{\VerbBar}{|}
\newcommand{\VERB}{\Verb[commandchars=\\\{\}]}
\DefineVerbatimEnvironment{Highlighting}{Verbatim}{commandchars=\\\{\}}
% Add ',fontsize=\small' for more characters per line
\usepackage{framed}
\definecolor{shadecolor}{RGB}{248,248,248}
\newenvironment{Shaded}{\begin{snugshade}}{\end{snugshade}}
\newcommand{\AlertTok}[1]{\textcolor[rgb]{0.94,0.16,0.16}{#1}}
\newcommand{\AnnotationTok}[1]{\textcolor[rgb]{0.56,0.35,0.01}{\textbf{\textit{#1}}}}
\newcommand{\AttributeTok}[1]{\textcolor[rgb]{0.13,0.29,0.53}{#1}}
\newcommand{\BaseNTok}[1]{\textcolor[rgb]{0.00,0.00,0.81}{#1}}
\newcommand{\BuiltInTok}[1]{#1}
\newcommand{\CharTok}[1]{\textcolor[rgb]{0.31,0.60,0.02}{#1}}
\newcommand{\CommentTok}[1]{\textcolor[rgb]{0.56,0.35,0.01}{\textit{#1}}}
\newcommand{\CommentVarTok}[1]{\textcolor[rgb]{0.56,0.35,0.01}{\textbf{\textit{#1}}}}
\newcommand{\ConstantTok}[1]{\textcolor[rgb]{0.56,0.35,0.01}{#1}}
\newcommand{\ControlFlowTok}[1]{\textcolor[rgb]{0.13,0.29,0.53}{\textbf{#1}}}
\newcommand{\DataTypeTok}[1]{\textcolor[rgb]{0.13,0.29,0.53}{#1}}
\newcommand{\DecValTok}[1]{\textcolor[rgb]{0.00,0.00,0.81}{#1}}
\newcommand{\DocumentationTok}[1]{\textcolor[rgb]{0.56,0.35,0.01}{\textbf{\textit{#1}}}}
\newcommand{\ErrorTok}[1]{\textcolor[rgb]{0.64,0.00,0.00}{\textbf{#1}}}
\newcommand{\ExtensionTok}[1]{#1}
\newcommand{\FloatTok}[1]{\textcolor[rgb]{0.00,0.00,0.81}{#1}}
\newcommand{\FunctionTok}[1]{\textcolor[rgb]{0.13,0.29,0.53}{\textbf{#1}}}
\newcommand{\ImportTok}[1]{#1}
\newcommand{\InformationTok}[1]{\textcolor[rgb]{0.56,0.35,0.01}{\textbf{\textit{#1}}}}
\newcommand{\KeywordTok}[1]{\textcolor[rgb]{0.13,0.29,0.53}{\textbf{#1}}}
\newcommand{\NormalTok}[1]{#1}
\newcommand{\OperatorTok}[1]{\textcolor[rgb]{0.81,0.36,0.00}{\textbf{#1}}}
\newcommand{\OtherTok}[1]{\textcolor[rgb]{0.56,0.35,0.01}{#1}}
\newcommand{\PreprocessorTok}[1]{\textcolor[rgb]{0.56,0.35,0.01}{\textit{#1}}}
\newcommand{\RegionMarkerTok}[1]{#1}
\newcommand{\SpecialCharTok}[1]{\textcolor[rgb]{0.81,0.36,0.00}{\textbf{#1}}}
\newcommand{\SpecialStringTok}[1]{\textcolor[rgb]{0.31,0.60,0.02}{#1}}
\newcommand{\StringTok}[1]{\textcolor[rgb]{0.31,0.60,0.02}{#1}}
\newcommand{\VariableTok}[1]{\textcolor[rgb]{0.00,0.00,0.00}{#1}}
\newcommand{\VerbatimStringTok}[1]{\textcolor[rgb]{0.31,0.60,0.02}{#1}}
\newcommand{\WarningTok}[1]{\textcolor[rgb]{0.56,0.35,0.01}{\textbf{\textit{#1}}}}
\usepackage{graphicx}
\makeatletter
\newsavebox\pandoc@box
\newcommand*\pandocbounded[1]{% scales image to fit in text height/width
  \sbox\pandoc@box{#1}%
  \Gscale@div\@tempa{\textheight}{\dimexpr\ht\pandoc@box+\dp\pandoc@box\relax}%
  \Gscale@div\@tempb{\linewidth}{\wd\pandoc@box}%
  \ifdim\@tempb\p@<\@tempa\p@\let\@tempa\@tempb\fi% select the smaller of both
  \ifdim\@tempa\p@<\p@\scalebox{\@tempa}{\usebox\pandoc@box}%
  \else\usebox{\pandoc@box}%
  \fi%
}
% Set default figure placement to htbp
\def\fps@figure{htbp}
\makeatother
\setlength{\emergencystretch}{3em} % prevent overfull lines
\providecommand{\tightlist}{%
  \setlength{\itemsep}{0pt}\setlength{\parskip}{0pt}}
\usepackage{bookmark}
\IfFileExists{xurl.sty}{\usepackage{xurl}}{} % add URL line breaks if available
\urlstyle{same}
\hypersetup{
  pdftitle={Assignment 2},
  pdfauthor={Bojan Endrovski},
  hidelinks,
  pdfcreator={LaTeX via pandoc}}

\title{Assignment 2}
\author{Bojan Endrovski}
\date{2026-04-20}

\begin{document}
\maketitle

\section{Assignment 2: Hypothesis testing and rejection
probabilities}\label{assignment-2-hypothesis-testing-and-rejection-probabilities}

\subsubsection{Bojan Endrovski}\label{bojan-endrovski}

\subsection{1. Test for one proportion}\label{test-for-one-proportion}

\subsubsection{Question 1}\label{question-1}

\begin{Shaded}
\begin{Highlighting}[]
\FunctionTok{prop.test}\NormalTok{(}\AttributeTok{x=}\DecValTok{13}\NormalTok{, }\AttributeTok{n=}\DecValTok{20}\NormalTok{, }\AttributeTok{p=}\FloatTok{0.5}\NormalTok{, }\AttributeTok{alternative=}\StringTok{"two.sided"}\NormalTok{)}
\end{Highlighting}
\end{Shaded}

\begin{verbatim}
## 
##  1-sample proportions test with continuity correction
## 
## data:  13 out of 20, null probability 0.5
## X-squared = 1.25, df = 1, p-value = 0.2636
## alternative hypothesis: true p is not equal to 0.5
## 95 percent confidence interval:
##  0.4094896 0.8369133
## sample estimates:
##    p 
## 0.65
\end{verbatim}

For the sample size of 20 we get a p-value of 0.2636 (\textgreater{}
0.05) thus failing to reject the null hypothesis. We don't have enough
evidence to claim that choosing the green door is significantly
different from 50\%.

\subsubsection{Question 2}\label{question-2}

\begin{Shaded}
\begin{Highlighting}[]
\FunctionTok{prop.test}\NormalTok{(}\AttributeTok{x=}\DecValTok{130}\NormalTok{, }\AttributeTok{n=}\DecValTok{200}\NormalTok{, }\AttributeTok{p=}\FloatTok{0.5}\NormalTok{, }\AttributeTok{alternative=}\StringTok{"two.sided"}\NormalTok{)}
\end{Highlighting}
\end{Shaded}

\begin{verbatim}
## 
##  1-sample proportions test with continuity correction
## 
## data:  130 out of 200, null probability 0.5
## X-squared = 17.405, df = 1, p-value = 3.02e-05
## alternative hypothesis: true p is not equal to 0.5
## 95 percent confidence interval:
##  0.5790769 0.7150579
## sample estimates:
##    p 
## 0.65
\end{verbatim}

For the sample size of 200 we get a p-value of 0.0000302 (\textless{}
0.05) thus rejecting the null hypothesis. We now have sufficient
evidence to claim that choosing the green door is significantly
different from 50\%. The proportion is the same (65\% of participants
chose the green door), but the bigger sample size give us more
statistical power.

\subsubsection{Question 3}\label{question-3}

\subsubsection{a)}\label{a}

\begin{Shaded}
\begin{Highlighting}[]
\NormalTok{n }\OtherTok{\textless{}{-}} \DecValTok{20}
\NormalTok{prob }\OtherTok{\textless{}{-}} \FloatTok{0.6}
\NormalTok{smp }\OtherTok{\textless{}{-}} \FunctionTok{rbinom}\NormalTok{(n, }\AttributeTok{size=}\DecValTok{1}\NormalTok{, }\AttributeTok{prob=}\NormalTok{prob)}
\FunctionTok{prop.test}\NormalTok{(}\FunctionTok{sum}\NormalTok{(smp), }\AttributeTok{n=}\NormalTok{n, }\AttributeTok{p=}\FloatTok{0.5}\NormalTok{, }\AttributeTok{alternative=}\StringTok{"two.sided"}\NormalTok{)}
\end{Highlighting}
\end{Shaded}

\begin{verbatim}
## 
##  1-sample proportions test with continuity correction
## 
## data:  sum(smp) out of n, null probability 0.5
## X-squared = 0.05, df = 1, p-value = 0.8231
## alternative hypothesis: true p is not equal to 0.5
## 95 percent confidence interval:
##  0.3204804 0.7617145
## sample estimates:
##    p 
## 0.55
\end{verbatim}

Running the experiment multiple times we get widely different results.
The sample estimates the probability anywhere from 0.3 to 0.7 and very
high p-value.

\subsubsection{b)}\label{b}

\begin{Shaded}
\begin{Highlighting}[]
\NormalTok{n }\OtherTok{\textless{}{-}} \DecValTok{200}
\NormalTok{prob }\OtherTok{\textless{}{-}} \FloatTok{0.6}
\NormalTok{smp }\OtherTok{\textless{}{-}} \FunctionTok{rbinom}\NormalTok{(n, }\AttributeTok{size=}\DecValTok{1}\NormalTok{, }\AttributeTok{prob=}\NormalTok{prob)}
\FunctionTok{prop.test}\NormalTok{(}\FunctionTok{sum}\NormalTok{(smp), }\AttributeTok{n=}\NormalTok{n, }\AttributeTok{p=}\FloatTok{0.5}\NormalTok{, }\AttributeTok{alternative=}\StringTok{"two.sided"}\NormalTok{)}
\end{Highlighting}
\end{Shaded}

\begin{verbatim}
## 
##  1-sample proportions test with continuity correction
## 
## data:  sum(smp) out of n, null probability 0.5
## X-squared = 4.205, df = 1, p-value = 0.0403
## alternative hypothesis: true p is not equal to 0.5
## 95 percent confidence interval:
##  0.5032062 0.6438648
## sample estimates:
##     p 
## 0.575
\end{verbatim}

In comparison with a), the estimates of the probability is a lot closer
to 0.6 but the p-value is under 0.05 on majority of the runs, rejecting
the null hypothesis. The closer the probability estimate is to the true
value of 0.6, the smaller the p-value.

\subsubsection{c)}\label{c}

\begin{Shaded}
\begin{Highlighting}[]
\NormalTok{n }\OtherTok{\textless{}{-}} \DecValTok{200}
\NormalTok{prob }\OtherTok{\textless{}{-}} \FloatTok{0.5}
\NormalTok{smp }\OtherTok{\textless{}{-}} \FunctionTok{rbinom}\NormalTok{(n, }\AttributeTok{size=}\DecValTok{1}\NormalTok{, }\AttributeTok{prob=}\NormalTok{prob)}
\FunctionTok{prop.test}\NormalTok{(}\FunctionTok{sum}\NormalTok{(smp), }\AttributeTok{n=}\NormalTok{n, }\AttributeTok{p=}\FloatTok{0.5}\NormalTok{, }\AttributeTok{alternative=}\StringTok{"two.sided"}\NormalTok{)}
\end{Highlighting}
\end{Shaded}

\begin{verbatim}
## 
##  1-sample proportions test with continuity correction
## 
## data:  sum(smp) out of n, null probability 0.5
## X-squared = 3.125, df = 1, p-value = 0.0771
## alternative hypothesis: true p is not equal to 0.5
## 95 percent confidence interval:
##  0.3657494 0.5067888
## sample estimates:
##     p 
## 0.435
\end{verbatim}

The p-value is high on most runs, not rejecting the null hypothesis even
when the probability estimate is very close to 0.5

\subsubsection{d)}\label{d}

\begin{Shaded}
\begin{Highlighting}[]
\NormalTok{n }\OtherTok{\textless{}{-}} \DecValTok{20}
\NormalTok{prob }\OtherTok{\textless{}{-}} \FloatTok{0.1}
\NormalTok{smp }\OtherTok{\textless{}{-}} \FunctionTok{rbinom}\NormalTok{(n, }\AttributeTok{size=}\DecValTok{1}\NormalTok{, }\AttributeTok{prob=}\NormalTok{prob)}
\FunctionTok{prop.test}\NormalTok{(}\FunctionTok{sum}\NormalTok{(smp), }\AttributeTok{n=}\NormalTok{n, }\AttributeTok{p=}\FloatTok{0.5}\NormalTok{, }\AttributeTok{alternative=}\StringTok{"two.sided"}\NormalTok{)}
\end{Highlighting}
\end{Shaded}

\begin{verbatim}
## 
##  1-sample proportions test with continuity correction
## 
## data:  sum(smp) out of n, null probability 0.5
## X-squared = 11.25, df = 1, p-value = 0.0007962
## alternative hypothesis: true p is not equal to 0.5
## 95 percent confidence interval:
##  0.01751264 0.33127760
## sample estimates:
##   p 
## 0.1
\end{verbatim}

Even with a small sample size we reject the null hypothesis on most
runs, unless the probability is estimated at avery different value than
0.1 (ex. prob=0.35 p-value=0.2636)

\subsubsection{Question 4}\label{question-4}

\begin{Shaded}
\begin{Highlighting}[]
\NormalTok{N }\OtherTok{\textless{}{-}} \DecValTok{500}
\NormalTok{n }\OtherTok{\textless{}{-}} \DecValTok{20}
\NormalTok{prob }\OtherTok{\textless{}{-}} \FloatTok{0.6}
\NormalTok{rej }\OtherTok{\textless{}{-}} \FunctionTok{rep}\NormalTok{(}\DecValTok{0}\NormalTok{, }\DecValTok{500}\NormalTok{)}
\ControlFlowTok{for}\NormalTok{ (i }\ControlFlowTok{in} \DecValTok{1}\SpecialCharTok{:}\NormalTok{N) \{}
\NormalTok{  smp }\OtherTok{\textless{}{-}} \FunctionTok{rbinom}\NormalTok{(n, }\AttributeTok{size=}\DecValTok{1}\NormalTok{, }\AttributeTok{prob=}\NormalTok{prob)}
\NormalTok{  res }\OtherTok{\textless{}{-}} \FunctionTok{prop.test}\NormalTok{(}\FunctionTok{sum}\NormalTok{(smp), }\AttributeTok{n=}\NormalTok{n, }\AttributeTok{p=}\FloatTok{0.5}\NormalTok{, }\AttributeTok{alternative=}\StringTok{"two.sided"}\NormalTok{)}
\NormalTok{  rej[i] }\OtherTok{\textless{}{-}}\NormalTok{ res}\SpecialCharTok{$}\NormalTok{p.value }\SpecialCharTok{\textless{}} \FloatTok{0.05}
\NormalTok{\}}
\FunctionTok{sprintf}\NormalTok{(}\StringTok{"a) Rejections: \%d out of \%d (probability: \%.4f)"}\NormalTok{, }\FunctionTok{sum}\NormalTok{(rej), N, }\FunctionTok{mean}\NormalTok{(rej))}
\end{Highlighting}
\end{Shaded}

\begin{verbatim}
## [1] "a) Rejections: 54 out of 500 (probability: 0.1080)"
\end{verbatim}

\begin{Shaded}
\begin{Highlighting}[]
\NormalTok{n }\OtherTok{\textless{}{-}} \DecValTok{200}
\NormalTok{prob }\OtherTok{\textless{}{-}} \FloatTok{0.6}
\NormalTok{rej }\OtherTok{\textless{}{-}} \FunctionTok{rep}\NormalTok{(}\DecValTok{0}\NormalTok{, }\DecValTok{500}\NormalTok{)}
\ControlFlowTok{for}\NormalTok{ (i }\ControlFlowTok{in} \DecValTok{1}\SpecialCharTok{:}\NormalTok{N) \{}
\NormalTok{  smp }\OtherTok{\textless{}{-}} \FunctionTok{rbinom}\NormalTok{(n, }\AttributeTok{size=}\DecValTok{1}\NormalTok{, }\AttributeTok{prob=}\NormalTok{prob)}
\NormalTok{  res }\OtherTok{\textless{}{-}} \FunctionTok{prop.test}\NormalTok{(}\FunctionTok{sum}\NormalTok{(smp), }\AttributeTok{n=}\NormalTok{n, }\AttributeTok{p=}\FloatTok{0.5}\NormalTok{, }\AttributeTok{alternative=}\StringTok{"two.sided"}\NormalTok{)}
\NormalTok{  rej[i] }\OtherTok{\textless{}{-}}\NormalTok{ res}\SpecialCharTok{$}\NormalTok{p.value }\SpecialCharTok{\textless{}} \FloatTok{0.05}
\NormalTok{\}}
\FunctionTok{sprintf}\NormalTok{(}\StringTok{"b) Rejections: \%d out of \%d (probability: \%.4f)"}\NormalTok{, }\FunctionTok{sum}\NormalTok{(rej), N, }\FunctionTok{mean}\NormalTok{(rej))}
\end{Highlighting}
\end{Shaded}

\begin{verbatim}
## [1] "b) Rejections: 402 out of 500 (probability: 0.8040)"
\end{verbatim}

\begin{Shaded}
\begin{Highlighting}[]
\NormalTok{n }\OtherTok{\textless{}{-}} \DecValTok{200}
\NormalTok{prob }\OtherTok{\textless{}{-}} \FloatTok{0.5}
\NormalTok{rej }\OtherTok{\textless{}{-}} \FunctionTok{rep}\NormalTok{(}\DecValTok{0}\NormalTok{, }\DecValTok{500}\NormalTok{)}
\ControlFlowTok{for}\NormalTok{ (i }\ControlFlowTok{in} \DecValTok{1}\SpecialCharTok{:}\NormalTok{N) \{}
\NormalTok{  smp }\OtherTok{\textless{}{-}} \FunctionTok{rbinom}\NormalTok{(n, }\AttributeTok{size=}\DecValTok{1}\NormalTok{, }\AttributeTok{prob=}\NormalTok{prob)}
\NormalTok{  res }\OtherTok{\textless{}{-}} \FunctionTok{prop.test}\NormalTok{(}\FunctionTok{sum}\NormalTok{(smp), }\AttributeTok{n=}\NormalTok{n, }\AttributeTok{p=}\FloatTok{0.5}\NormalTok{, }\AttributeTok{alternative=}\StringTok{"two.sided"}\NormalTok{)}
\NormalTok{  rej[i] }\OtherTok{\textless{}{-}}\NormalTok{ res}\SpecialCharTok{$}\NormalTok{p.value }\SpecialCharTok{\textless{}} \FloatTok{0.05}
\NormalTok{\}}
\FunctionTok{sprintf}\NormalTok{(}\StringTok{"c) Rejections: \%d out of \%d (probability: \%.4f)"}\NormalTok{, }\FunctionTok{sum}\NormalTok{(rej), N, }\FunctionTok{mean}\NormalTok{(rej))}
\end{Highlighting}
\end{Shaded}

\begin{verbatim}
## [1] "c) Rejections: 19 out of 500 (probability: 0.0380)"
\end{verbatim}

\begin{Shaded}
\begin{Highlighting}[]
\NormalTok{n }\OtherTok{\textless{}{-}} \DecValTok{20}
\NormalTok{prob }\OtherTok{\textless{}{-}} \FloatTok{0.1}
\NormalTok{rej }\OtherTok{\textless{}{-}} \FunctionTok{rep}\NormalTok{(}\DecValTok{0}\NormalTok{, }\DecValTok{500}\NormalTok{)}
\ControlFlowTok{for}\NormalTok{ (i }\ControlFlowTok{in} \DecValTok{1}\SpecialCharTok{:}\NormalTok{N) \{}
\NormalTok{  smp }\OtherTok{\textless{}{-}} \FunctionTok{rbinom}\NormalTok{(n, }\AttributeTok{size=}\DecValTok{1}\NormalTok{, }\AttributeTok{prob=}\NormalTok{prob)}
\NormalTok{  res }\OtherTok{\textless{}{-}} \FunctionTok{prop.test}\NormalTok{(}\FunctionTok{sum}\NormalTok{(smp), }\AttributeTok{n=}\NormalTok{n, }\AttributeTok{p=}\FloatTok{0.5}\NormalTok{, }\AttributeTok{alternative=}\StringTok{"two.sided"}\NormalTok{)}
\NormalTok{  rej[i] }\OtherTok{\textless{}{-}}\NormalTok{ res}\SpecialCharTok{$}\NormalTok{p.value }\SpecialCharTok{\textless{}} \FloatTok{0.05}
\NormalTok{\}}
\FunctionTok{sprintf}\NormalTok{(}\StringTok{"d) Rejections: \%d out of \%d (probability: \%.4f)"}\NormalTok{, }\FunctionTok{sum}\NormalTok{(rej), N, }\FunctionTok{mean}\NormalTok{(rej))}
\end{Highlighting}
\end{Shaded}

\begin{verbatim}
## [1] "d) Rejections: 492 out of 500 (probability: 0.9840)"
\end{verbatim}

\begin{enumerate}
\def\labelenumi{\alph{enumi})}
\tightlist
\item
  The probability of rejection is low due to the small sample
\item
  The probability of rejection is is much higher because of the bigger
  sample
\item
  The probability of rejection is low because the null hypothesis is
  indeed true
\item
  The probability of rejection is high, in-spite of the small sample due
  to the strong statistical effect
\end{enumerate}

\subsection{2. Test for two proportions}\label{test-for-two-proportions}

\subsubsection{Question 1}\label{question-1-1}

\begin{Shaded}
\begin{Highlighting}[]
\NormalTok{xs }\OtherTok{\textless{}{-}} \FunctionTok{c}\NormalTok{(}\DecValTok{18}\NormalTok{, }\DecValTok{28}\NormalTok{)}
\NormalTok{ns }\OtherTok{\textless{}{-}} \FunctionTok{c}\NormalTok{(}\DecValTok{30}\NormalTok{, }\DecValTok{40}\NormalTok{)}
\FunctionTok{prop.test}\NormalTok{(}\AttributeTok{x=}\NormalTok{xs, }\AttributeTok{n=}\NormalTok{ns, }\AttributeTok{alternative=}\StringTok{"two.sided"}\NormalTok{)}
\end{Highlighting}
\end{Shaded}

\begin{verbatim}
## 
##  2-sample test for equality of proportions with continuity correction
## 
## data:  xs out of ns
## X-squared = 0.38176, df = 1, p-value = 0.5367
## alternative hypothesis: two.sided
## 95 percent confidence interval:
##  -0.3547755  0.1547755
## sample estimates:
## prop 1 prop 2 
##    0.6    0.7
\end{verbatim}

From this experiment we can't reject the null hypothesis that px == py.
The is no sufficient evidence that the social interaction influences
door choice.

\subsubsection{Question 2}\label{question-2-1}

\begin{Shaded}
\begin{Highlighting}[]
\NormalTok{sz1 }\OtherTok{\textless{}{-}}\DecValTok{30}
\NormalTok{sz2 }\OtherTok{\textless{}{-}} \DecValTok{40}
\NormalTok{smp1 }\OtherTok{\textless{}{-}} \FunctionTok{rbinom}\NormalTok{(sz1, }\DecValTok{1}\NormalTok{, }\FloatTok{0.6}\NormalTok{)}
\NormalTok{smp2 }\OtherTok{\textless{}{-}} \FunctionTok{rbinom}\NormalTok{(sz2, }\DecValTok{1}\NormalTok{, }\FloatTok{0.8}\NormalTok{)}
\NormalTok{xs }\OtherTok{\textless{}{-}} \FunctionTok{c}\NormalTok{(}\FunctionTok{sum}\NormalTok{(smp1), }\FunctionTok{sum}\NormalTok{(smp2))}
\NormalTok{ns }\OtherTok{\textless{}{-}} \FunctionTok{c}\NormalTok{(sz1, sz2)}
\FunctionTok{prop.test}\NormalTok{(}\AttributeTok{x=}\NormalTok{xs, }\AttributeTok{n=}\NormalTok{ns, }\AttributeTok{alternative=}\StringTok{"two.sided"}\NormalTok{)}
\end{Highlighting}
\end{Shaded}

\begin{verbatim}
## 
##  2-sample test for equality of proportions with continuity correction
## 
## data:  xs out of ns
## X-squared = 0.625, df = 1, p-value = 0.4292
## alternative hypothesis: two.sided
## 95 percent confidence interval:
##  -0.3643341  0.1310007
## sample estimates:
##    prop 1    prop 2 
## 0.6333333 0.7500000
\end{verbatim}

The null hypothesis is not rejected on more often that it is rejected,
but with the sample sizes of 30 and 40 is hard to derive a conclusive
answer. When the probability estimates are further apart, the p-value is
smaller.

\end{document}
